% GENERATED FILE. Edit canonical lessons, then run npm run generate:books.
% canonical-source-hash: fnv1a64:f804c40bdaef6d52
% canonical-lessons: ES-C603-privilegio, ES-C603-igualdad, ES-C603-desigualdad, ES-C603-justicia, ES-C603-injusto

\chapter{Derechos}
\label{ch:es-derechos}
\hlchaptermodality{603}

\begin{chapteropening}
\textbf{By the end of this chapter:} \emph{I can talk about equality, privilege, justice and what is unfair.}

Complete the last lesson of chapter 603: I can talk about equality, privilege, justice and what is unfair.
\end{chapteropening}

\section[el privilegio]{el privilegio — the privilege}

\label{lesson:ES-C603-privilegio}



\begin{quote}
Before the new one: say the Spanish for the voter, then the Spanish for the democracy.
\end{quote}

\subsection*{You'll want to know: el privilegio}
\textbf{el privilegio} — \textquotedblleft{}the privilege\textquotedblright{}.

\begin{quote}
Estudiar en esta universidad es un privilegio.
\end{quote}

\begin{quote}
\emph{Studying at this university is a privilege.}
\end{quote}

\emph{Privilegiado} — privileged.

\begin{cousinweb}
\emph{privilegio} is a learned borrowing from Latin \emph{privilegium} 'a law made for one person' (EtymDB), from \emph{privus} 'one's own' and \emph{lex} 'law'.
\end{cousinweb}

\begin{grammarlens}[title={What you've built}]
One more word for telling, explaining and describing at B1.
\end{grammarlens}

\subsection*{Your turn}
\begin{itemize}
\raggedright
  \item \emph{Say it:} \emph{el privilegio}
  \item \emph{Say it:} \emph{el privilegio}, once more
  \item \emph{Say it:} a sentence of your own with \emph{el privilegio}
  \item \emph{Recall:} say the Spanish for the voter, then the Spanish for the democracy, then say \emph{el privilegio} again
\end{itemize}

\begin{tcolorbox}[breakable,colback=teal!4,colframe=teal!35!black,title={Before you move on}]
What does el privilegio mean? (\textquotedblleft{}The privilege\textquotedblright{}.) Say it once more.
\end{tcolorbox}
\section[la igualdad]{la igualdad — the equality}

\label{lesson:ES-C603-igualdad}



\begin{quote}
Before the new one: say the Spanish for the democracy, then the Spanish for the privilege.
\end{quote}

\subsection*{You'll want to know: la igualdad}
\textbf{la igualdad} — \textquotedblleft{}the equality\textquotedblright{}.

\begin{quote}
Luchamos por la igualdad.
\end{quote}

\begin{quote}
\emph{We fight for equality.}
\end{quote}

\emph{Igual} — equal, the same.

\begin{cousinweb}
\emph{igualdad} is inherited from Latin \emph{aequalitas} 'evenness, equality', built on \emph{aequus} 'level, equal' (EtymDB).
\end{cousinweb}

\begin{grammarlens}[title={What you've built}]
One more word for telling, explaining and describing at B1.
\end{grammarlens}

\subsection*{Your turn}
\begin{itemize}
\raggedright
  \item \emph{Say it:} \emph{la igualdad}
  \item \emph{Say it:} \emph{la igualdad}, once more
  \item \emph{Say it:} a sentence of your own with \emph{la igualdad}
  \item \emph{Recall:} say the Spanish for the democracy, then the Spanish for the privilege, then say \emph{la igualdad} again
\end{itemize}

\begin{tcolorbox}[breakable,colback=teal!4,colframe=teal!35!black,title={Before you move on}]
What does la igualdad mean? (\textquotedblleft{}The equality\textquotedblright{}.) Say it once more.
\end{tcolorbox}
\section[la desigualdad]{la desigualdad — the inequality}

\label{lesson:ES-C603-desigualdad}



\begin{quote}
Before the new one: say the Spanish for the privilege, then the Spanish for the equality.
\end{quote}

\subsection*{You'll want to know: la desigualdad}
\textbf{la desigualdad} — \textquotedblleft{}the inequality\textquotedblright{}.

\begin{quote}
La desigualdad crece en las ciudades.
\end{quote}

\begin{quote}
\emph{Inequality is growing in the cities.}
\end{quote}

The opposite of \emph{la igualdad}.

\begin{cousinweb}
\emph{desigualdad} is built on \emph{igualdad} with the prefix \emph{des-}; \emph{igualdad} is inherited from Latin \emph{aequalitas}, from \emph{aequus} 'equal' (EtymDB).
\end{cousinweb}

\begin{grammarlens}[title={What you've built}]
One more word for telling, explaining and describing at B1.
\end{grammarlens}

\subsection*{Your turn}
\begin{itemize}
\raggedright
  \item \emph{Say it:} \emph{la desigualdad}
  \item \emph{Say it:} \emph{la desigualdad}, once more
  \item \emph{Say it:} a sentence of your own with \emph{la desigualdad}
  \item \emph{Recall:} say the Spanish for the privilege, then the Spanish for the equality, then say \emph{la desigualdad} again
\end{itemize}

\begin{tcolorbox}[breakable,colback=teal!4,colframe=teal!35!black,title={Before you move on}]
What does la desigualdad mean? (\textquotedblleft{}The inequality\textquotedblright{}.) Say it once more.
\end{tcolorbox}
\section[la justicia]{la justicia — the justice}

\label{lesson:ES-C603-justicia}



\begin{quote}
Before the new one: say the Spanish for the equality, then the Spanish for the inequality.
\end{quote}

\subsection*{You'll want to know: la justicia}
\textbf{la justicia} — \textquotedblleft{}the justice\textquotedblright{}.

\begin{quote}
Pedimos justicia.
\end{quote}

\begin{quote}
\emph{We ask for justice.}
\end{quote}

Also the courts: \emph{acudir a la justicia} — to go to court.

\begin{cousinweb}
\emph{justicia} is a learned borrowing from Latin \emph{iustitia} 'justice', built on \emph{iustus} 'lawful, right' (EtymDB).
\end{cousinweb}

\begin{grammarlens}[title={What you've built}]
One more word for telling, explaining and describing at B1.
\end{grammarlens}

\subsection*{Your turn}
\begin{itemize}
\raggedright
  \item \emph{Say it:} \emph{la justicia}
  \item \emph{Say it:} \emph{la justicia}, once more
  \item \emph{Say it:} a sentence of your own with \emph{la justicia}
  \item \emph{Recall:} say the Spanish for the equality, then the Spanish for the inequality, then say \emph{la justicia} again
\end{itemize}

\begin{tcolorbox}[breakable,colback=teal!4,colframe=teal!35!black,title={Before you move on}]
What does la justicia mean? (\textquotedblleft{}The justice\textquotedblright{}.) Say it once more.
\end{tcolorbox}
\section[injusto]{injusto — unfair}

\label{lesson:ES-C603-injusto}



\begin{quote}
Before the new one: say the Spanish for the inequality, then the Spanish for the justice.
\end{quote}

\subsection*{You'll want to know: injusto}
\textbf{injusto} — \textquotedblleft{}unfair\textquotedblright{}.

\begin{quote}
Es injusto pagar más.
\end{quote}

\begin{quote}
\emph{It is unfair to pay more.}
\end{quote}

\emph{Es injusto que} + subjunctive.

\begin{cousinweb}
\emph{injusto} is a learned borrowing from Latin \emph{iniustus} 'unfair', built from \emph{in-} 'not' and \emph{iustus} 'right' (Lewis \& Short).
\end{cousinweb}

\begin{grammarlens}[title={What you've built}]
One more word for telling, explaining and describing at B1.
\end{grammarlens}

\subsection*{Your turn}
\begin{itemize}
\raggedright
  \item \emph{Say it:} \emph{injusto}
  \item \emph{Say it:} \emph{injusto}, once more
  \item \emph{Say it:} a sentence of your own with \emph{injusto}
  \item \emph{Recall:} say the Spanish for the inequality, then the Spanish for the justice, then say \emph{injusto} again
\end{itemize}

\begin{tcolorbox}[breakable,colback=teal!4,colframe=teal!35!black,title={Before you move on}]
What does injusto mean? (\textquotedblleft{}Unfair\textquotedblright{}.) Say it once more.
\end{tcolorbox}
